\section{Parameters, ranks, and probability scales}
\label{sec:parameters}

This Appendix collects the choices made throughout the proof and
checks their dependencies. All quantities chosen before \(n\) are
constants. They can be very large; no upper bound on their size is
needed, and none of these choices depends on an adversarial unit law.
The fixed mixers may be chosen separately at each sufficiently large
\(n\), by \Cref{lem:generic-mixers}.

\subsection{The acyclic order of construction}

The early constants in
\eqref{eq:early-constants}--\eqref{eq:parameter-order} are
\begin{align}
 C_0&=1000,& g&=10^9+1,&
 D&=4C_0g,& M&=2^{1000},\nonumber\\
 k_{\max}&=56(g-1),&
 \zeta&=\frac{1}{1000k_{\max}},&
 K_1&=\left\lceil\frac{4(D+1)}{\zeta}\right\rceil .
 \label{eq:ledger-early}
\end{align}
In particular \(g\) is odd and \(D=4000g\).
The early rank budgets are
\begin{align}
 r&=2K_1+2,&
 L_0&=\lceil100(D+10)\rceil,&
 d_0&=4rL_0,\nonumber\\
 u_0&=10(K_1+d_0+1),&
 K&\in\N,\qquad
 K\ge\frac{10(D+K_1+u_0+10)}{\zeta}.
 \label{eq:ledger-law}
\end{align}
These are the bounds in
\eqref{eq:law-budgets} and \eqref{eq:pin-budgets};
the construction takes \(K\) to be the least integer meeting
the displayed lower bound.  Only the first peeling is used.
The auxiliary number \(u_0\) is retained in the numerical definition
of a conservative bound \(K\ge K_1\); it does not count another
operation in the present construction.

The remaining early tolerances in
\Cref{thm:phase-alternative,lem:injection} include
\begin{equation}
 c_s=10^{-12}(1+r)^{-4},\qquad
 q_0=2^{-100r-10},\qquad
 \delta_{\mathrm{cov}}=.0005.
 \label{eq:ledger-cover-tolerances}
\end{equation}
For each \(n\), the moment order is the largest power of two
\(s\le c_sN\); it is even for sufficiently large \(n\), and
\(s\ge c_sN/2\).
Choose \(p_0>0\) sufficiently small in terms of \(r,D,c_s,q_0\)
and the fixed phase tolerance. Fix the transpose-injection error
thresholds as well, with their total over all component and endpoint
roles below the required constant error, such as \(.001\).
These thresholds precede selector and channel sizes.
For the bounded-tuple transpose tests in
\Cref{lem:channel-thinning}, we may fix
\[
 q_{\rm test}=d_0+L_0+2r+2K+10.
\]
This covers the tiny covers, fixed fresh probes, and bounded pin
and marked-profile tests.  The growing lists of length proportional
to \(N\) in \Cref{lem:injection} instead use an absolute
reference-law count, as explained in its proof.

The exponential demands of the no-cover argument are simultaneously
satisfiable in this order. Fix a phase-test threshold \(\vartheta>0\)
smaller than the required fixed phase accuracy, and choose a constant
\(\Gamma>D+2r+1\) large enough to absorb the indicated density and
bounded-rank tensor counts. With \(L_\vartheta=\log_2(1/\vartheta)\),
it suffices to arrange
\begin{align}
 \frac{c_s}{2}
 \left(\frac{q_0}{2}\log_2(1/p_0)-2-L_\vartheta\right)
     &>\Gamma,\label{eq:ledger-p0-demand}\\
 \frac{c_s}{2}
 \left(\frac{hq_0}{4}-L_\vartheta\right)
     &>\Gamma.\label{eq:ledger-h-moment-demand}
\end{align}
The first determines a sufficiently small \(p_0\) using only early
quantities; the second is a later lower bound on \(h\).
Indeed the two moment terms are
\(2^{2s}p_0^{q_0s/2}\) and \(2^{-hq_0s/4}\), and Markov at threshold
\(\vartheta\) costs \(2^{sL_\vartheta}\).
The comparison between raw and independent channel laws has a factor
depending on \(h\), but it is constant in \(n\) and does not alter the
linear exponents.

For the bounded tiny-cover record use
\[
 R_0=2^{100(r+1)^2(d_0+|\cE|+1)},\qquad
 \delta_{\rm ph}=\frac{p_0\delta_{\mathrm{cov}}}{64R_0},\qquad
 \delta_* =10^{-6}\delta_{\rm ph}/64.
\]
The preparation argument supplies these early retained-mass bounds
and \(c_*=\delta_*^2/64>0\), as in
\eqref{eq:preparation-early-margins}.  They measure retained unit
mass and successful pair mass in the original pooled law with
endpoint bound \(M\pi\).  Choose
\begin{equation}\label{eq:ledger-colours}
 0<\eta_{\mathrm{col}}<c_*/4,\qquad
 0<a<\min\{1/100,\eta_{\mathrm{col}}/800\}.
\end{equation}
A colour partition is refined into leaves with their actual weights;
preparation is performed on their pooled law.  Thus these margins
are independent of both the number of colours and the smallest
positive colour mass.  Indeed original equal-colour pair mass is
at most \(\sum_c p_c^2\le\eta_{\mathrm{col}}\); after restriction
of mass \(\delta_*\) its normalized mass is at most
\(\eta_{\mathrm{col}}/\delta_*^2\).  Either calculation gives
strictly positive successful mass after the colour exclusion.

Now prescribe \(0<\beta\le1\), and choose \(L>10^6/\beta\), as
in \eqref{eq:transfer-a-L}.  Larger values of \(\beta\) follow
by weakening the graph assertions.  For the matrix obstruction
one chooses \(\beta\) sufficiently small in terms of \(a\),
which is permitted here.  Fix the additional bounded marginal
factors and positive injection/phase tolerances required by the
asymmetric test with this \(L\).  These later fixed tolerances
may depend on \(L\); they do not change the preceding
\(\delta_*,c_*,\eta_{\mathrm{col}}\).  In particular the two
preparation problems need not have the same success margin.

Next choose the row and selector parameters:
\begin{align}
 R_*&=2(K+20),&
 j_*&\ge10(K+1)(R_*+20),\nonumber\\
 B_0&=2|\cE|K(R_*+14),&
 B_*&=10(B_0+10000g+1),\nonumber\\
 2^{b-2j_*}&>100g(B_*+1),&
 p_*&=\sum_{j=0}^{j_*}\binom bj,\nonumber\\
 A&=10(K+1)p_*(g^2+5),&
 s_0&=4\lceil A\rceil,\qquad
 J>100(s_0+K^2+1).
 \label{eq:ledger-selectors}
\end{align}
Integer inequalities are met by rounding upward when needed.
This agrees with \Cref{lem:label-support,lem:generic-mixers}.
In particular \(j_*\) exceeds the interpolation degree required for
\((K+1)(R_*+14)\) labels. The length \(b\) is chosen before \(A\),
and hence before \(s_0,J\).

Now choose a sufficiently large integer \(r_0\), always setting
\begin{equation}
                         h=1000r_0.
 \label{eq:ledger-channels}
\end{equation}
The demands are finite in number and depend only on preceding
constants:
\begin{enumerate}
\item the moment demand \eqref{eq:ledger-h-moment-demand};
\item fixed-cover and fresh-probe transpose-injection errors in
  \Cref{thm:phase-alternative};
\item the unconditional transpose-injection estimate in \Cref{lem:injection};
\item residual ranks and private-channel space in
  \Cref{thm:gradient-realization};
\item the tester rank inequality in \Cref{lem:zero-delta}.
\end{enumerate}
Each is eventually satisfied by increasing \(r_0\).
Quantitative checks appear below.

After fixing \(h\), choose the thinning rate \(\rho>0\).
It is enough to require
\begin{equation}\label{eq:ledger-thinning-rate}
 10^4|\cE|2^{2h}\rho
   <\min\{\zeta/100,10^{-6},\epsilon_{\mathrm{img}},
                  \epsilon_{\mathrm{inj}}\},
\end{equation}
where the positive constants \(\epsilon_{\mathrm{img}}\) and
\(\epsilon_{\mathrm{inj}}\) are chosen below the reserved
full-frame image and injection rate margins.  Those margins
are fixed before this choice.  Put
\[
 \kappa_{\mathrm{thin}}=2|\cE|(2^h-1),\qquad
 \tau_{\mathrm{thin}}=10|\cE|2^h\rho.
\]
The thinning lemma gives
\begin{equation}\label{eq:ledger-thinned-density}
 \pi^2\le2^{\tau_{\mathrm{thin}}N}\rawlaw^2,
 \qquad 3000g+\tau_{\mathrm{thin}}<D-1.
\end{equation}
Thus the distribution-test pair cap leaves a fixed positive
reserve below the absolute cap \(2^{DN}\rawlaw^2\), even after
bounded normalizations and removal of exponentially light colours.
The survival fraction is asymptotic to
\(2^{-\kappa_{\mathrm{thin}}\rho N}\).
Because \(\kappa_{\mathrm{thin}}\rho<.01\), the simultaneous
additive concentration error \(2^{-.1N}\) in the construction
of the thinning sets is exponentially smaller than this fraction.
Its use in conditional probabilities is therefore justified.

Finally choose a sufficiently large integer \(M_0\), also relative
to \(1/\rho\), and put
\begin{equation}
                  N=M_0n,\qquad m=2^{C_0gN}.
 \label{eq:ledger-final-dimensions}
\end{equation}
The coefficient of every \(O(n)\) count paid as a fraction of \(N\)
has then been fixed. Choose \(M_0\) to make all those fractions
smaller than their prescribed slacks. Only afterward does \(n\)
tend to infinity through integers with \(n\ge2r_0\).
The resulting dependency order is
\[
 \begin{gathered}
 (C_0,g,D,M,k_{\max},\zeta,K_1)
 \longrightarrow
 (r,L_0,d_0,u_0,K,\text{early tolerances},p_0,\delta_*,c_*)\\
 \longrightarrow(\eta_{\mathrm{col}},a)
 \longrightarrow(\beta,L,\text{asymmetric tolerances})\\
 \longrightarrow(R_*,j_*,B_0,B_*,b,p_*,A,s_0,J)
 \longrightarrow(r_0,h)\longrightarrow\rho
 \longrightarrow M_0\longrightarrow n .
 \end{gathered}
\]
Here \(b\) is only the selector length; \(\beta\) is the graph
threshold.  The fixed palette has size at most
\(\lceil1/(2a)\rceil+1\), and changes only bounded complexities
and the sufficiently-large-\(n\) onset.

\subsection{The channel demands precede later tests}

In the frequent-output argument, a fixed collection of \(L_0\)
fresh probes succeeds on B-mass at least
\begin{equation}
       q_*=(\delta_{\mathrm{cov}}/2)
              (\delta_{\mathrm{cov}}/4)^{L_0}>0.
 \label{eq:ledger-probe-mass}
\end{equation}
The one-vertex density factor depends only on the early marked-law
restrictions and \(p_0\).
Transpose-injection failure on the fixed fresh vectors is at most
that factor times \(O(L_0)2^{L_0-h}\).
It can be made smaller than \(q_*/2\) by increasing \(h\), without
consulting a selector choice or a future common test.
The later raw-event exponent is at most
\[
       -(.99-.51)L_0N+DN+\text{small linear errors}.
\]
The specified \(L_0\) leaves a margin close to \((.48L_0-D)N\).
The small coefficient-count errors are paid by \(M_0\).

For \Cref{lem:injection} only the unconditional accepting
family is needed.  The independent primal-span part uses the
unchanged primal marginals.  For the channel part, the
estimate with \(\ell=\lfloor N/(10(K+1))\rfloor\) compares
\(c^\ell\) with \(2^{-(h-2K)\ell+o(N)}\), where \(c>0\)
is fixed by the required failure tolerance.  A sufficient demand is
\begin{equation}\label{eq:ledger-injection-demand}
 h>2K+\log_2(1/c)+30(K+1).
\end{equation}
It leaves a positive linear exponent margin.  The thinning loss
in \eqref{eq:ledger-thinning-rate} is chosen smaller than that
margin; fixed marginal factors cost only constants.  In particular
this estimate does not require a table-dependent tolerance or a
new pinning step.

For channel realization, the explicit bounds in
\Cref{eq:baseline-rank-bound,eq:residual-rank-bound,eq:block-compression-bound,eq:quotient-compression-bound}
are
\[
\begin{aligned}
 B_{\mathrm{lin}}&=3K+28,&
 R_{\mathrm{res}}&=14K+140,\\
 D_{\mathrm{blk}}&=2p_*|\cE|(K+14+R_{\mathrm{res}}),&
 D_{\mathrm{quo}}&=2p_*\bigl(1+(g^2+3)D_{\mathrm{blk}}\bigr)+2K.
\end{aligned}
\]
The pure correction rank per component is at most
\(30r_0+28J|\cE|+D_{\mathrm{quo}}\), and the remaining channel
pairing has rank at least \(h-2K-4B_{\mathrm{lin}}\).
Thus \Cref{eq:realization-r0-requirement} asks precisely for
\begin{equation}
       970r_0>28J|\cE|+D_{\mathrm{quo}}+2K+4B_{\mathrm{lin}}.
 \label{eq:ledger-realization-demand}
\end{equation}
Every quantity on the right is independent of \(h,r_0,n\)
and is fixed by parameters through \(J\).
The baseline first factors through bounded pin/key projections on
primal inputs, then derivative corrections of bounded rank are
chosen, and the small residual is compressed before the
\(r_0\)-sized pure terms are factored through private channels.

The tester inequality in \Cref{lem:zero-delta} is
\begin{equation}
 \frac{g-1}{2}\left(g-\frac{4D}{g}\right)r_0
          >2(g-1)(h+2JK_1)+1.
 \label{eq:ledger-pure-tester-demand}
\end{equation}
Substituting \(D=4000g\) and \(h=1000r_0\) gives
\[
       \frac{g-20000}{2}\,r_0
             >4JK_1+\frac1{g-1}.
\]
The coefficient on the left is positive for the specified \(g\).
This is therefore another compatible lower bound on \(r_0\).

\subsection{Pin, atom, and tensor-rank ledger}

\begin{center}
\begin{tabular}{@{}p{.39\textwidth}p{.53\textwidth}@{}}
\toprule
Quantity & Bound and role\\
\midrule
First pin dimension & \(K_1\), summed over components and signs.\\
Pin dimension used throughout & At most \(K_1\), hence at most
the conservative budget \(K\).\\
Tiny-cover dimension & At most \(d_0=4rL_0\); its offset
record has a bounded number of possibilities.\\
Effective space & \(\dim C_i\le K^2\); every member has total component
rank at most \(K\).\\
Attack atoms at one endpoint & At most \(14\), across its two lists.\\
Component rank in the annihilator argument &
\(2(K+14)\le R_*=2(K+20)\).\\
List length per cross edge & At most \(7\) on either side.\\
Atom positions in one unit's full query & At most \(4\cdot7=28\).\\
General key-vector slots &
\(28\cdot2(g-1)=k_{\max}=56(g-1)\).\\
\bottomrule
\end{tabular}
\end{center}

Each tag atom occupies \(g-1\) star components and supplies one plus
and one minus vector at each. These are vector counts, not scalar
Gram counts.
The label exclusions in \Cref{lem:sparse-pin-labels} give nominal
independence modulo the pins, as proved in
\Cref{lem:small-tables}, part~\textup{(ii)}.
The effective-space dimension is counted in projected pin bases:
for a fixed endpoint \(\sum_e\dim S_{i,e}^\pm\le K\), so
\[
 \sum_e\dim(S_{i,e}^+\otimes S_{i,e}^-)
 \le
 \left(\sum_e\dim S_{i,e}^+\right)
 \left(\sum_e\dim S_{i,e}^-\right)
 \le K^2.
\]
Thus an effective tensor basis is encoded in bounded projected
coordinates, not as unrestricted matrices with \(O(n^2)\) entries.

The unprojected representation of one gradient in
\Cref{thm:gradient-realization} has per-component rank at most
\[
                         15r_0+14J|\cE|.
\]
Two endpoint contributions give \(30r_0+O(1)\).
If \(A_+,A_-\) remove bounded pin/key subspaces, then
\[
 \rank(M-A_+\trans M A_-)
       \le\rank(I-A_+)+\rank(I-A_-).
\]
Hence this projection contributes only bounded residual rank,
without multiplying the tester rank by a pin-dependent constant.
The later blockwise compression and derivative corrections preserve
that independence from \(r_0\).

\subsection{Scalar records and the final dimension multiplier}

Write
\begin{equation}
                  d_n=\dim\cB=p_*\bigl(1+(g^2+3)n\bigr).
 \label{eq:ledger-column-dimension}
\end{equation}
Then
\[
       \frac{d_n}{N}
       =\frac{p_*(g^2+3)}{M_0}+\frac{p_*}{M_0n}.
\]
This ratio, and the ratio after adding the bounded channel dimension,
can be made arbitrarily small by choosing \(M_0\) last.
In particular the image estimate can use a fixed sufficiently small
\(\varepsilon<\zeta/4\).  For the tiny-cover calculation impose,
in addition, that the total primal coefficient-count and image
loss be smaller than \(\rho/12\).  The same choice leaves room
for the injective frames and prescribed self-Gram values.

The collision records are linear in \(n\).
There are at most \(56\) point inputs across two units, costing at most
\[
                         56(g^2+3)n+56b
\]
bits. A list of \(K\) nominal pin directions uses at most
\(2K(d_n+h)\) coefficient bits, plus bounded component, sign and rank
indices.
Once pin projections are fixed, a basis of \(C_i\) uses at most
\(K^4\) binary coefficients in its tensor-coordinate space;
there are four endpoints.
Tester summaries, numerical small tables, projection relationships
and finite statuses have bounded lengths in \(n\).
The construction matrices \(E,L,R\) are fixed data and are not
encoded afresh in each query record.

The total frozen dimension of a unit is at most \(K+k_{\max}\),
over all components and signs.
There are \(2(d_n+h)\) nominal columns of a given sign at each
component. Recording all interactions of both units' columns with
opposite frozen images therefore uses at most
\[
                         4(d_n+h)(K+k_{\max})
\]
scalar bits.
The deterministic solution is selected from these data and the point
parameters. No unrestricted \(O(n^2)\)-entry matrix is appended.
The number of record cells, including bounded nominal choices when
needed, is consequently at most
\begin{equation}
                         2^{C_{\mathrm{rec}}n+C_{\mathrm{rec},0}},
 \label{eq:ledger-record-count}
\end{equation}
for constants fixed before \(M_0\).

Only mixed primal-channel Gram entries are tested in the final target.
For one cross Gram block and component, the two endpoint primal
summands have total dimension \(2d_n\) and the channel summands have
total dimension \(2h\) on either side. The two mixed products
contain at most \(8d_nh\) entries.
There are two cross Gram blocks; hence at most
\begin{equation}
                              16|\cE|\,h\,d_n
 \label{eq:ledger-tested-gram-bits}
\end{equation}
bits are tested. Frozen or redundant entries only reduce this count.
This is \(O(n)\), with coefficient already fixed.

Choose \(M_0\) to meet every earlier coefficient/codimension slack
and, in particular, so that for sufficiently large \(n\),
\begin{equation}
 \log_2(\text{record-cell count})<.005N,\qquad
 \text{tested scalar Gram bits}<.01N.
 \label{eq:ledger-collision-bit-slacks}
\end{equation}
The smaller image and pin-count slacks form a finite collection
and can be met by the same final enlargement of \(M_0\).

The conditional point-mass estimate on a surviving cell is
\[
 \begin{split}
 \log_2 p_{\max}/N
  &\le -(1-2\zeta)(k+t)+(k+.01)\\
  &=-(1-2\zeta)t+.01+2\zeta k
       \le-.95t
       \qquad(t\ge1,\ k\le k_{\max}).
 \end{split}
\]
Here \(2\zeta k\le.002\).
A nonzero quotient character of rank \(t\) has mean at most
\[
       2^{tN/2}
       \bigl(2^{-.95tN}2^{-.95tN}\bigr)^{1/2}
             =2^{-.45tN}.
\]
This pays for fewer than \(.01N\) scalar bits, leaving agreement
probability at least \(2^{-.02N}\).
For each fixed parameter pair and key, pruning cells below
\(2^{-(k+.01)N}\) loses overlap at most
\[
 2|\cK_n|\,2^{C_{\mathrm{rec}}n+C_{\mathrm{rec},0}}
       2^{-(k+.01)N},
\]
which is exponentially small since \(|\cK_n|\le2^{kN}\) and
\eqref{eq:ledger-collision-bit-slacks} holds.
This uses Fubini at a fixed parameter pair, without intersecting
pruning decisions over all opposite parameters.

\subsection{Probability scales and later analytic choices}
\label{sec:probability-scales}

The marginal distinctions used in the proof are summarized below.
A bounded multiplier may depend on a fixed preparation restriction;
it is independent of \(n\).

\begin{center}
\begin{tabular}{@{}p{.36\textwidth}p{.56\textwidth}@{}}
\toprule
Mechanism & Scale and permissible cost\\
\midrule
Thinned reference law & Primal marginal exactly that of
\(\rawlaw\); full pair density at most
\(2^{\tau_{\mathrm{thin}}N}\rawlaw^2\).\\
Original tested unit law & Joint cap \(Q_N\pi^2\); both endpoint
marginals at most \(M\pi\), or \(L\pi\) for the asymmetric test.\\
Prepared mixed law & Absolute joint cap within the reserved
\(2^{DN}\rawlaw^2\) budget; bounded marginals relative to
\(\pi\), hence bounded primal marginals relative to the original
primal law.\\
Normalized leaves & A common \(2^{C_{\mathrm{leaf}}N}\rawlaw^2\)
cap and fresh-image min-entropy with at most \(2\zeta\) slack;
no bounded individual full-frame marginal is asserted.\\
Fixed-cost preparation restrictions & Positive mass fixed before
\(h\); bounded logarithmic costs are absorbed by rate reserves.\\
Single-endpoint typicality & Error \(2^{-\Omega(n)}\), tested on
primal frames; arbitrary channel-dependent flags remain flags.\\
Paired oversampling and cross-rank failures & Error
\(2^{-\Omega(n^2)}\), surviving \(2^{O(N)}=2^{O(n)}\)
joint inflation.\\
Tiny-cover compatible pairs & Positive mass fixed before \(h\),
using either the identically zero full-sum phase or the additional
thinning gain.\\
Key equality & Factor \(|\cK_n|^{-1}\ge2^{-kN}\),
\(k\le k_{\max}\).\\
Final four-hole probability &
\(2^{-(k_{\max}+.03)N-o(N)}\), exceeding
\(2^{-100gN}\) for sufficiently large \(n\).\\
\bottomrule
\end{tabular}
\end{center}

For clarity, the extra gain from thinning is not obtained by
subtracting the full-frame density loss from \(\rho\).
On a fixed tiny-cover record, the companion tuple of length
\(t\) has primal point probability at most
\(2^{-tN+\rho N/12}\).  If the full output is nonzero, one
endpoint must have a nonzero channel combination in a prescribed
nonzero translate of its allowed set.  Conditional on its primal
frame, \Cref{lem:channel-thinning} bounds that event by
\(2^{-\rho N/2}\).  The point probability of companion plus
nonzero full output is therefore at most
\(2^{-tN-\rho N/3}\), after absorbing fixed normalizations.
The Walsh estimate pairs one such bound with a companion-only
bound, giving at most
\[
 2^{(t_A+t_B)N/2}
 \bigl(2^{-t_AN+\rho N/12}
       2^{-t_BN-\rho N/3}\bigr)^{1/2}
 \le2^{-\rho N/12}.
\]
There is no exponential union over these records: one of their
boundedly many values was fixed at an early constant mass cost.
If the full output is zero on the retained record, the lifted
phase vanishes for every ordered retained pair.  Both cases
supply constant compatibility mass.  No further peeling or
inverse-polynomial compatibility branch is used.

Fixed-cost restrictions are compatible with the leaf trimming in
\Cref{lem:peeling}.  More generally, if a later analytical restriction
has total mass \(q_n=2^{-o(N)}\), leaves on which its relative mass
is below \(2^{-\zeta N}\) contribute at most
\[
 2^{-\zeta N}/q_n=2^{-(\zeta-o(1))N}
\]
of the restricted law.  On the other leaves division costs at most
\(2^{\zeta N}\), changing \((1-\zeta)tN\) to at worst
\((1-2\zeta)tN\) for every positive rank \(t\).  This trimming
preserves the old pins.  Marginal hypotheses for later empirical
common tests concern the primal projection; none requires replacing
\(\pi\) by \(\rawlaw\) in a full-frame marginal inequality.

Partitions, entropy bounds, weak-product accuracies and finite
transformation sample sizes are analytical choices made after the
construction constants.  They may depend on \(M_0\), a fixed positive
test-cell mass, or a fixed accuracy.  They do not change pin ranks,
key-slot counts or the scalar record coefficient in
\eqref{eq:ledger-record-count}.  For each finite common-test stage
these choices change only constants and the required onset in
\(n\); a diagonal subsequence passes each fixed finite stage in
turn.  No rate uniform over all such stages is needed.

For example, unary positivity uses
\(\gamma=2^{-K^2-3}\) and \(b_{\mathrm{int}}=2^{-J+2s_0}<\gamma\).
On a reference key cell of mass at least \(\delta>0\), the proof
first obtains conditional input mass at least \(\delta/2\), so it
can use
\[
 u>\frac{2}{\delta(\gamma-b_{\mathrm{int}})^2}
\]
transformations.  This changes the constant in
\(O_u(2^{-s_0n})\), while \(s_0=4\lceil A\rceil\) and the profile
count \(2^{An}\) remain fixed.  Thus a smaller cell does not require
a larger construction parameter \(J\).

The final collision margin is
\[
 100g-[56(g-1)+.03]=44g+55.97>0.
\]
This positive exponent gap absorbs any fixed positive subsequential
overlap constant.  A violating sequence for either distribution
test would yield one limit model with positive successful overlap,
which contradicts its asserted upper bound \(2^{-100gN}\).
Uniformity over all limit models is unnecessary.

Finally, \Cref{sec:graph-transfer} uses
\(Q_N=2^{3000gN}\) and \(m=2^{1000gN}\).  Entropy fingerprints
have logarithmic count
\[
 O(N^3 2^{100gN})=o(m).
\]
For either fixed positive linear choice of \(k\), the coloured
exceptional-pair estimate is
\[
 m^{2k}(4/Q_N)^k=2^{-(1000gN-2)k}.
\]
The coloured position exception has mass less than \(8/M\);
at \(k=\lceil m/100\rceil\) its tail is at most
\((800e/M)^k\).  Fewer than \(4/\eta_{\mathrm{col}}\)
exceptional colours, each of size at most \(2am\), account for
fewer than \(m/100\) matching edges by
\eqref{eq:ledger-colours}.  In the asymmetric transfer use
\(k=\lceil\beta m/100\rceil\); vertex and point exceptional
sets have mass below \(1/L\), with tail at most
\((100e/(\beta L))^k\).  These bounds pay for the
\(\exp(o(m))\) terminal union and prove the graph properties
at \(m/20\) and \(\beta m/10\), respectively.
